\section{Simultaneous labeling for discounted monotone maps}
\label{sec:labeling}

We now construct the labels that will permit a box to shrink while
retaining the integer comparison vectors around the discounted fixed
point.  The construction depends only on order and on the response to
a common translation.  We therefore work in the following abstract
setting: $n\ge1$ and $F:\R^n\to\R^n$ satisfy
\begin{equation}
 t\le t'\ \Longrightarrow\ F(t)\le F(t'),\qquad
 F(t+c\one)=F(t)+\gamma c\one
 \quad(t\in\R^n,\ c\in\R),\qquad 0<\gamma<1.
 \label{eq:label-map}
\end{equation}
All vector inequalities, minima, maxima, and floors in this section
are coordinatewise.  The algorithm uses $F$ only by evaluating
$\lfloor F(t)\rfloor$ at integer vectors $t$.

The two pivot passes and subsequent mass change adapt the simultaneous
labeling construction of
\cite[Sections~4--5]{OpenAI2026Deterministic}.  Here the operator has
the discounted translation rule in Equation~\eqref{eq:label-map}.
The proof below establishes the required labeling theorem directly;
the signs of the witness translations are what make that adaptation
possible.

\subsection{Witnesses and the labeling guarantee}

An input consists of $A\in\Z^n$, a width
\begin{equation}
 D=h2^d,\qquad 128\le h\le255,\quad h,d\in\Z,\quad d\ge0,
 \label{eq:label-width}
\end{equation}
and positive rational mass vectors $a,b\in\mathbb Q_{>0}^n$.
Write $u=A+D\one$, $P=D/32$, and
$a(S)=\sum_{i\in S}a_i$, $b(S)=\sum_{i\in S}b_i$ for
$S\subseteq\{1,\ldots,n\}$.  The representation in
Equation~\eqref{eq:label-width} is unique.  Its short range of
mantissas will later allow small relative decreases in the width.

\begin{definition}[Witnesses and claims]
\label{def:label-witness}
An integer vector $y\in[A,u]$ is a \emph{subsolution in the box} if
\[
 F_i(y)\ge y_i\quad\text{for every }i\in S_y,
 \qquad S_y=\{i:y_i>A_i\}.
\]
It is a \emph{plus witness} if also $a(S_y)\le1$; its claims are
the indices $i$ for which $y_i\ge u_i-P$.

An integer vector $z\in[A,u]$ is a \emph{supersolution in the box} if
\[
 F_i(z)\le z_i\quad\text{for every }i\in S_z,
 \qquad S_z=\{i:z_i<u_i\}.
\]
It is a \emph{minus witness} if also $b(S_z)\le1$; its claims are
the indices $i$ for which $z_i\le A_i+P$.
The sets $S_y,S_z$ are the respective supports.  A witness is
\emph{relevant} if it has at least one claim.
\end{definition}

Every claim belongs to its witness's support, since $0<P<D$.
The output must respect all these claims simultaneously, although
the algorithm is not given any witness and does not search for one.

\begin{theorem}[Recursive simultaneous labeling]
\label{thm:labeling}
Let $F:\R^n\to\R^n$ satisfy Equation~\eqref{eq:label-map}.
There is a terminating deterministic procedure $\cL(A,D,a,b)$,
using exact evaluations of $\lfloor F(t)\rfloor$ for $t\in\Z^n$
and integer and rational arithmetic, with the following guarantee.
For every input as in Equation~\eqref{eq:label-width} and
Definition~\ref{def:label-witness}, its output
$\ell\in\{+,-\}^n$ satisfies
\[
 \begin{aligned}
 y_i\ge A_i+D-P&\ \Longrightarrow\ \ell_i=+
       &&\text{for every plus witness }y,\\
 z_i\le A_i+P&\ \Longrightarrow\ \ell_i=-
       &&\text{for every minus witness }z.
 \end{aligned}
\]
\end{theorem}

The existence of a labeling will already follow from the finite box
iteration below.  The purpose of the recursive construction is to
obtain such labels with the controlled branching and decreasing rank
of Lemma~\ref{lem:label-rank}, replacing dependence on the numerical
width by dependence on its exponent.

\begin{lemma}[Comparison and a finite box iteration]
\label{lem:label-comparison}
Every subsolution $y$ and supersolution $z$ in the same box satisfy
$y\le z$.  Starting at $t^0=u$, the iteration
\begin{equation}
 t^{m+1}=T(t^m),\qquad
 T(t)=\min\{u,\max\{A,\lfloor F(t)\rfloor\}\},
 \label{eq:label-floor}
\end{equation}
stops at a vector $t^*$ after at most $nD+1$ evaluations of $T$.
For every such pair $y,z$ one has $y\le t^*\le z$.
\end{lemma}

\begin{proof}
If $m=\max_j(y_j-z_j)>0$, choose $i$ with $y_i-z_i=m$.
Since $y_i>z_i\ge A_i$ and $z_i<y_i\le u_i$, both defining
inequalities apply at $i$.  Monotonicity and the common translation
rule give
\[
 m\le F_i(y)-F_i(z)
   \le F_i(z+m\one)-F_i(z)=\gamma m,
\]
a contradiction.

The map $T$ preserves the finite integer box and is monotone.
Moreover $T(u)\le u$, so induction gives $t^{m+1}\le t^m$.
Each changing step decreases the nonnegative integer
$\sum_i(t_i^m-A_i)$ by at least one.  Its initial value is $nD$,
which proves the termination bound, including the final unchanged
step.

For an integer $y\in[A,u]$, the subsolution condition is equivalent
to $y\le T(y)$.  Indeed, when $y_i>A_i$, the latter inequality is
equivalent to $y_i\le\lfloor F_i(y)\rfloor$, and at the other
coordinates it is automatic.  If $y$ is a subsolution, induction
gives $y\le t^m$: from $y\le t^m$ one obtains
$y\le T(y)\le T(t^m)=t^{m+1}$.  Thus $y\le t^*$.
The fixed point $t^*=T(t^*)$ is itself a subsolution, and the
comparison principle gives $t^*\le z$ for every supersolution $z$.
\end{proof}

\subsection{The recursive procedure}

A \emph{centered call} of even width $r$ at an integer pivot $k$
means the call $\cL(k-(r/2)\one,r,a',b')$ with the indicated
masses.  All calls use the same map $F$.  In a nonbase call set
\begin{equation}
 c=A+\frac D2\one,\qquad M=\frac D8,\qquad
 q=\frac D{128},\qquad p=\frac q{32}=\frac D{4096}.
 \label{eq:label-scales}
\end{equation}
The pivot will stay between $c-M\one$ and $c+M\one$.
The two passes use narrow centered boxes to identify some labels
directly and arrange the remaining witnesses so that the later mass
change preserves their admissibility. The downward pass preserves
the plus masses, and the upward pass preserves the minus masses.
A final wider centered call supplies the signs used to change both
mass vectors.

\medskip
\noindent\textbf{Procedure $\cL(A,D,a,b)$.}
Execute the first applicable base case.  If neither applies, execute
Steps~3--6 in order.
\begin{enumerate}
\item \textbf{Mass base case.}
If $a_i b_i>1$ for every $i$, return $+$ at $i$ when $b_i>1$,
and return $-$ otherwise.

\item \textbf{Width base case.}
If $d<12$, compute $t^*$ by Equation~\eqref{eq:label-floor}.
Return $+$ at $i$ when $t_i^*>A_i+D/2$, and return $-$ otherwise.

\item \textbf{Downward pass.}
Initialize $k=c$.  Make a centered width-$q$ call at $k$ with
masses $(a,\frac43 b)$, and use its labels to update all coordinates
simultaneously by
\begin{equation}
 k_i^{\mathrm{new}}=
 \begin{cases}
  \max\{c_i-M,k_i-p\},&\text{if the child labels }i\text{ by }-,\\
  k_i,&\text{if the child labels }i\text{ by }+.
 \end{cases}
 \label{eq:label-down-update}
\end{equation}
Replace $k$ by $k^{\mathrm{new}}$ and repeat until an update leaves
the pivot unchanged.  Denote this last pivot by $k^{\downarrow}$ and set
\begin{equation}
 I_- =\{i:k_i^{\downarrow}<c_i-P\}.
 \label{eq:label-minus-override}
\end{equation}

\item \textbf{Upward pass.}
Start at $k=k^{\downarrow}$.  Make a centered width-$q$ call at $k$
with masses $(\frac43 a,b)$, and update simultaneously by
\begin{equation}
 k_i^{\mathrm{new}}=
 \begin{cases}
  \min\{c_i+M,k_i+p\},&\text{if the child labels }i\text{ by }+,\\
  k_i,&\text{if the child labels }i\text{ by }-.
 \end{cases}
 \label{eq:label-up-update}
\end{equation}
Replace $k$ by $k^{\mathrm{new}}$ and repeat until an unchanged
update.  Denote the last pivot by $k^f$ and set
\begin{equation}
 I_+ =\{i:k_i^f>c_i+P\}.
 \label{eq:label-plus-override}
\end{equation}

\item \textbf{Provisional labels and new masses.}
Make a centered width-$D/2$ call at $k^f$ with masses $(a,b)$.
For its output $\ell^0$, define
\begin{equation}
 (\bar a_i,\bar b_i)=
 \begin{cases}
  (\frac45 a_i,\frac85 b_i),&\ell_i^0=+,\\[1mm]
  (\frac85 a_i,\frac45 b_i),&\ell_i^0=-.
 \end{cases}
 \label{eq:label-reweight}
\end{equation}

\item \textbf{Continuation and overrides.}
Call $\cL(A,D,\bar a,\bar b)$.  In its output replace the signs
on $I_-$ by $-$ and those on $I_+\setminus I_-$ by $+$.
Return this label vector.
\end{enumerate}

Each pass includes the final child call whose update changes nothing.
That call will certify a mass inequality at the stopped pivot.  All
coordinates in one update use the same old pivot and the same child
output; there is no order-dependent coordinate update.

The purpose of the passes is the following mass property. For every
witness with a claim outside its own override set---$I_+$ for a plus
witness and $I_-$ for a minus witness---the coordinates in its support
whose provisional labels match its sign will have total mass greater
than $3/4$. Here mass means $a$-mass for a plus witness and $b$-mass
for a minus witness. Its original total mass is at most one, so
Step~5 will leave its new total mass below
\[
                  \frac85-\frac45\cdot\frac34=1.
\]
It will therefore remain a witness for the continuation in the
unchanged box. At the same time, every product $a_i b_i$ is multiplied
by $32/25$. We first use this product growth to prove termination;
the subsequent witness analysis proves the stated mass property and
hence correctness.

\begin{lemma}[Valid boxes and bounded passes]
\label{lem:label-mechanics}
Every child input is valid and its box lies inside its parent's box.
If the child calls terminate, each pivot pass makes at most
$1024n+1$ calls.  The width base case makes at most
$522240n+1$ evaluations of $\lfloor F\rfloor$.
\end{lemma}

\begin{proof}
In a nonbase call $d\ge12$, so $c,M,q,p$ and the half-widths
$q/2,D/4$ are integral.  The centered widths have the allowed
representations $q=h2^{d-7}$ and $D/2=h2^{d-1}$.
The pivot caps keep $c-M\one\le k\le c+M\one$, and a centered
box has width $r\le D/2$.  Its boundaries satisfy
\[
 k-\frac r2\one\ge A+\frac D8\one\ge A,
 \qquad
 k+\frac r2\one\le A+\frac{7D}8\one\le u.
\]
The continuation uses the same box.  All masses remain positive
rationals.

Every pivot coordinate lies on the grid of spacing $p$ through
$c_i$, because $M/p=512$ is integral.  In a changing update some
coordinate moves by $p$; each coordinate is monotone during a pass
and travels at most $2M$, hence moves at most $1024$ times.
There are at most $1024n$ changing updates, followed by one
unchanged update.  Finally, $d<12$ implies
$D\le255\cdot2^{11}=522240$, so the width bound follows from
Lemma~\ref{lem:label-comparison}.
\end{proof}

The continuation preserves the width, so width reduction alone does
not prove termination.  The mass products provide the missing
measure.  Define
\begin{equation}
 \mu(a,b)=\min_i a_i b_i,\qquad \rho=\frac{32}{25},\qquad
 B(\mu)=\min\{j\in\Z_{\ge0}:\rho^j\mu>1\}.
 \label{eq:label-budget}
\end{equation}
This integer is used only in the proof, not computed by the procedure.

\begin{lemma}[Rank and branching]
\label{lem:label-rank}
The procedure terminates.  At a nonbase call every child has smaller
rank $d+B(\mu(a,b))$.  There is exactly one child preserving $B$,
the provisional call, and it decreases $d$ by one.  There are at most
$2048n+3$ other children; each decreases $B$ by at least one and
does not increase $d$.
\end{lemma}

\begin{proof}
At a nonbase call $B(\mu)>0$, since otherwise the mass base case
would apply.  Whenever $\mu'\ge\rho\mu$, the definition gives
$B(\mu')\le B(\mu)-1$.  Both pivot passes multiply every mass
product by $4/3\ge\rho$.  The continuation multiplies every mass
product by $(4/5)(8/5)=\rho$, independently of the provisional
label.  These children reduce $B$ and preserve or reduce $d$.
Only the provisional call keeps the masses; it halves the width
and reduces $d$ by one.  This proves the rank assertion.

Induction on the nonnegative rank proves termination.  The base cases
are finite by their definitions and Lemma~\ref{lem:label-mechanics}.
At a nonbase call every child terminates by induction, and the bound
on each pivot pass then applies.  The number of children decreasing
$B$ is at most $2(1024n+1)+1=2048n+3$.
\end{proof}

\subsection{Translating a witness into a centered box}

To analyze the labels, we follow an arbitrary parent witness through
the centered calls.  Translation places its largest deviation from
the pivot at one boundary of the child box; clipping removes
coordinates far below that largest deviation.  The surviving support
will identify the mass forced near the largest deviation when a
pivot pass stops.

Fix a nonbase input and a pivot $k\in[c-M\one,c+M\one]$.
For a relevant plus witness $y$, define its \emph{peak} and
\emph{gaps} by
\begin{equation}
 s_y(k)=\max_i(y_i-k_i),\qquad
 g_i^y(k)=s_y(k)-(y_i-k_i).
 \label{eq:label-plus-gaps}
\end{equation}
For a relevant minus witness $z$, put
\begin{equation}
 s_z(k)=\max_i(k_i-z_i),\qquad
 g_i^z(k)=s_z(k)-(k_i-z_i).
 \label{eq:label-minus-gaps}
\end{equation}
The gaps are nonnegative integers and at least one gap is zero.
The following lemma handles both signs, including the inequalities
introduced by the discount factor.

\begin{lemma}[Translation and clipping]
\label{lem:label-translation}
Let $r\in\{q,D/2\}$ and let $y$ or $z$ be a relevant witness
of the parent box.  Define, according to its sign,
\begin{equation}
 \begin{aligned}
 y'_i&=\max\{k_i-r/2,\ y_i-s_y(k)+r/2\},\\
 z'_i&=\min\{k_i+r/2,\ z_i+s_z(k)-r/2\}.
 \end{aligned}
 \label{eq:label-translation}
\end{equation}
The resulting vector is an integer subsolution or supersolution,
respectively, in the centered width-$r$ box.  Its support is exactly
$\{i:g_i(k)<r\}$, which is contained in the parent's support.
It claims every index with $g_i(k)\le r/32$.  It is a witness
for that child whenever its support has mass at most one for the
child's mass vector of the corresponding sign.
\end{lemma}

\begin{proof}
A parent claim gives, in either case,
\begin{equation}
 s(k)\ge D/2-M-P=11D/32.
 \label{eq:label-peak-lower}
\end{equation}
Outside the parent support one has $y_i=A_i$ or $z_i=u_i$.
The corresponding deviation is at most $-D/2+M=-3D/8$,
so its gap is at least $23D/32>D/2$.  Such an index cannot
belong to $\{i:g_i(k)<r\}$.

Before clipping, the two expressions on the right of
Equation~\eqref{eq:label-translation} are
\[
 k_i+r/2-g_i^y(k),\qquad
 k_i-r/2+g_i^z(k),
\]
respectively.  They lie on the correct side of the far boundary;
clipping at the other boundary puts the vector in the child box.
These expressions also show directly that the support is exactly
the set of gaps strictly below $r$, and that every gap at most
$r/32$ gives a claim.

For the plus case let $t=-s_y(k)+r/2<0$, where the strict sign
follows from Equation~\eqref{eq:label-peak-lower} and $r/2\le D/4$.
One has $y'\ge y+t\one$.  At an index in the child support,
equality holds in that coordinate, and the parent subsolution
inequality applies.  Therefore
\[
 F_i(y')\ge F_i(y+t\one)=F_i(y)+\gamma t
          \ge y_i+t=y'_i,
\]
using $\gamma t\ge t$.  For the minus case let
$t=s_z(k)-r/2>0$.  Now $z'\le z+t\one$, with coordinatewise
equality on the child support, and
\[
 F_i(z')\le F_i(z+t\one)=F_i(z)+\gamma t
          \le z_i+t=z'_i,
\]
using $\gamma t\le t$.  Thus the required inequalities survive
both translation and clipping.  The stated mass condition is
exactly the remaining condition for being a child witness.
\end{proof}

\subsection{What the two pivot passes enforce}

We prove correctness by induction on the rank from
Lemma~\ref{lem:label-rank}.  For the remainder of this subsection,
fix a nonbase call and assume that all its children satisfy
Theorem~\ref{thm:labeling}.  Witnesses, supports, and claims refer
to the parent input unless stated otherwise.

In each pass, the sign whose masses are unchanged protects its own
peak: coordinates close enough to that peak must receive its sign
and remain stationary.  The opposite sign has its masses multiplied
by $4/3$.  At the final unchanged update, any witness of that sign
with a claim outside its override set must consequently have mass greater than
$3/4$ near its peak.  The next two lemmas make these statements
precise.

\begin{lemma}[The downward pass]
\label{lem:label-down}
For every relevant plus witness $y$, the peak $s_y(k)$ is constant
throughout the downward pass and at most $D/2$.  No plus witness
claims an index in $I_-$.  Every minus witness $z$ with a claim
outside $I_-$ satisfies
\begin{equation}
 b\bigl(\{i:g_i^z(k^{\downarrow})<q\}\bigr)>3/4.
 \label{eq:label-down-mass}
\end{equation}
\end{lemma}

\begin{proof}
Fix a relevant plus witness $y$ and consider an update from $k$.
The width-$q$ translation in Lemma~\ref{lem:label-translation}
is a plus witness for the child, since the plus masses are unchanged
and its support is contained in $S_y$.  Its claims include all
indices with $g_i^y(k)\le q/32=p$.  The child labels all of them
by $+$, so they do not move.  In particular, every peak achiever
continues to attain the old peak.

At any other coordinate let $\delta_i=k_i-k_i^{\mathrm{new}}$.
Then $0\le\delta_i\le p$ and
\[
 y_i-k_i^{\mathrm{new}}
   =s_y(k)-g_i^y(k)+\delta_i<s_y(k).
\]
Thus no moving coordinate reaches the old peak, even when several
coordinates move simultaneously.  The peak is unchanged.  At the
initial pivot $c$, $y\le u$ gives $s_y(c)\le D/2$.  If $y$
claims $i$, it follows that
\[
 k_i^{\downarrow}\ge y_i-D/2
    \ge A_i+D/2-P=c_i-P,
\]
so $i\notin I_-$.

For the mass assertion, let $z$ claim $j\notin I_-$.  Then
$k_j^{\downarrow}\ge c_j-P$ and $z_j\le A_j+P$, whence
\begin{equation}
 s_z(k^{\downarrow})\ge D/2-2P=7D/16>D/2-M.
 \label{eq:label-down-uncapped}
\end{equation}
At a lower-capped coordinate $k_i^{\downarrow}=c_i-M$ one instead
has $k_i^{\downarrow}-z_i\le D/2-M$, since $z_i\ge A_i$.
Every peak achiever is therefore strictly above the lower cap.

If the mass in Equation~\eqref{eq:label-down-mass} were at most
$3/4$, the width-$q$ translation of $z$ would have support mass
at most one under the final child's minus masses $\frac43 b$.
It would be a minus witness and would claim every peak achiever.
The child would label such a coordinate by $-$, forcing it to move
down because it is above the lower cap.  This contradicts the
unchanged update.  The mass is therefore strictly greater than
$3/4$.
\end{proof}

\begin{lemma}[The upward pass]
\label{lem:label-up}
For every relevant minus witness $z$, the peak $s_z(k)$ is constant
throughout the upward pass and at most $D/2$.  No minus witness
claims an index in $I_+$.  At the final pivot the following mass
inequalities hold:
\begin{align}
 b\bigl(\{i:g_i^z(k^f)<q\}\bigr)&>3/4
 &&\text{if $z$ has a claim outside $I_-$},
 \label{eq:label-final-minus-mass}\\
 a\bigl(\{i:g_i^y(k^f)<q\}\bigr)&>3/4
 &&\text{if $y$ has a claim outside $I_+$}.
 \label{eq:label-final-plus-mass}
\end{align}
\end{lemma}

\begin{proof}
Fix a relevant minus witness $z$.  The width-$q$ translation is a
minus witness for every child in this pass because the minus masses
are unchanged.  All coordinates with $g_i^z(k)\le p$ must receive
$-$ and stay fixed.  In particular a peak achiever stays fixed.
At a coordinate with $g_i^z(k)>p$, write
$\delta_i=k_i^{\mathrm{new}}-k_i\in[0,p]$.  Then
\[
 k_i^{\mathrm{new}}-z_i
    =s_z(k)-g_i^z(k)+\delta_i<s_z(k).
\]
No coordinate overtakes the peak, so the peak is constant.  At the
starting pivot $k^{\downarrow}\le c$ and $z\ge A$, it is at most
$D/2$.  A minus claim at $i$ consequently implies
\[
 k_i^f\le z_i+D/2\le A_i+P+D/2=c_i+P,
\]
excluding $i$ from $I_+$.  Because the peak is constant and $k$
only increases during this pass, every gap $g_i^z(k)$ is
nonincreasing.  Thus the set of gaps below $q$ can only grow,
and Equation~\eqref{eq:label-down-mass} implies
Equation~\eqref{eq:label-final-minus-mass}.

Now let a plus witness $y$ claim $j\notin I_+$.  At the final
pivot $k_j^f\le c_j+P$ and $y_j\ge A_j+D-P$, so
\[
 s_y(k^f)\ge D/2-2P=7D/16>D/2-M.
\]
An upper-capped coordinate $k_i^f=c_i+M$ has
$y_i-k_i^f\le D/2-M$, since $y_i\le u_i$.  Every peak
achiever is strictly below the upper cap.  If the mass in
Equation~\eqref{eq:label-final-plus-mass} were at most $3/4$,
its width-$q$ translation would be a plus witness for the last
child, whose plus masses are $\frac43 a$.  Every peak achiever
would have label $+$ and would move up, contradicting the final
unchanged update.  This proves the remaining inequality.
\end{proof}

The two passes have now separated the claims into those that can
safely receive an override and those whose witnesses have substantial
mass near a peak at $k^f$.  The wider provisional call labels all
that mass with the witness's sign.  Reducing the masses of that sign
then keeps each remaining witness admissible for the continuation.
Figure~\ref{fig:labeling-passes} summarizes how the two passes supply
the information used in this final step.

\begin{figure}[!htb]
 \centering
 \begin{tikzpicture}[
  font=\small,
  stage/.style={draw=black!60,rounded corners=2pt,align=left,
    text width=5.7cm,inner sep=8pt},
  flow/.style={-{Stealth[length=2mm]},thick,draw=black!65}
]
\node[stage,fill=blue!4] (down) at (0,0) {
  \textbf{Downward pass: $(a,\frac43 b)$}\\[3pt]
  Plus peaks stay fixed.\\
  $I_-$ has no plus claims.\\[3pt]
  At stopping, each minus witness\\
  claiming outside $I_-$ has\\
  $b$-mass $>3/4$ at gaps $<q$.
};
\node[stage,fill=orange!6,right=9mm of down] (up) {
  \textbf{Upward pass: $(\frac43 a,b)$}\\[3pt]
  Minus peaks stay fixed.\\
  $I_+$ has no minus claims.\\[3pt]
  At stopping, each plus witness\\
  claiming outside $I_+$ has\\
  $a$-mass $>3/4$ at gaps $<q$.
};
\draw[flow] (down.east) -- (up.west);
\node[align=center,font=\small,below=6mm of up] (persist)
  {The earlier minus-mass bound persists:\\
   every minus gap is nonincreasing.};
\draw[flow] (up.south) -- (persist.north);
\node[stage,text width=10.5cm,align=center,fill=black!3,
  below=27mm of $(down.south)!0.5!(up.south)$] (wide) {
  \textbf{Provisional call: width $D/2$, masses $(a,b)$}\\[3pt]
  Every gap $<q$ is a child claim, since $q<D/64$.\\
  Each remaining witness has mass $>3/4$\\
  receiving its sign.
};
\draw[flow] (down.south) -- ++(0,-1.28) -| (wide.north west);
\draw[flow] (persist.south) -- ++(0,-.15) -| (wide.north east);
\node[align=center,below=6mm of wide] (last) {
  Multiply masses at matching provisional signs by $4/5$,\\
  and at opposite signs by $8/5$.  The witness's total mass becomes $<1$.
};
\draw[flow] (wide.south) -- (last.north);
\end{tikzpicture}

 \caption{The two passes preserve opposite peaks and force opposite
 mass inequalities.  A remaining witness means one with a claim
 outside its own override set: $I_+$ for plus witnesses and $I_-$
 for minus witnesses.  The upward pass preserves the mass inequality
 obtained during the downward pass.}
 \label{fig:labeling-passes}
\end{figure}

\subsection{Mass reweighting and completion of the proof}

\begin{proof}[Proof of Theorem~\ref{thm:labeling}]
Termination follows from Lemma~\ref{lem:label-rank}; we prove the
labeling guarantee by induction on the same rank.  In the mass
base case, an index labeled $+$ has $b_i>1$, so no minus witness
can contain that index in its support or claim it.  At an index
labeled $-$ one has $b_i\le1$ and $a_i b_i>1$, hence $a_i>1$;
no plus witness can claim that index.  This proves both guarantees.

In the width base case, Lemma~\ref{lem:label-comparison} gives
$y\le t^*\le z$ for all witnesses.  A plus claim forces
$t_i^*\ge u_i-P>A_i+D/2$; a minus claim forces
$t_i^*\le A_i+P<A_i+D/2$.  Thus the returned signs are correct.

For a nonbase call, assume correctness of all its children.  Let
$Q_+=\{i:\ell_i^0=+\}$ and $Q_-=\{i:\ell_i^0=-\}$ be the
provisional sign sets.  Consider a plus witness $y$ that has a
claim outside $I_+$.  By Equation~\eqref{eq:label-final-plus-mass},
the set $G_y=\{i:g_i^y(k^f)<q\}$ has $a$-mass greater than
$3/4$.  It lies in $S_y$ by Lemma~\ref{lem:label-translation}.
Apply that lemma with width $r=D/2$ to the provisional call.
Its plus masses are unchanged, so the translated vector is a
plus witness.  Since
\[
 q=D/128<D/64=r/32,
\]
every index in $G_y$ is a claim of this translated witness.
All these indices belong to $Q_+$ by the child's guarantee.
Consequently
\begin{equation}
 \bar a(S_y)
   =\frac85 a(S_y)-\frac45 a(S_y\cap Q_+)
   <\frac85-\frac45\cdot\frac34=1.
 \label{eq:label-plus-survives}
\end{equation}
The parent vector $y$ is still a subsolution in the unchanged box,
so it is a plus witness for the continuation.

For a minus witness $z$ having a claim outside $I_-$, the set
$G_z=\{i:g_i^z(k^f)<q\}$ has $b$-mass greater than $3/4$ by
Equation~\eqref{eq:label-final-minus-mass}.  It lies in $S_z$.
The width-$D/2$ minus translation is a witness for the provisional
call because the minus masses are unchanged; every index in $G_z$
is one of its claims, and hence belongs to $Q_-$.  Therefore
\begin{equation}
 \bar b(S_z)
   =\frac85 b(S_z)-\frac45 b(S_z\cap Q_-)
   <\frac85-\frac45\cdot\frac34=1.
 \label{eq:label-minus-survives}
\end{equation}
The same vector $z$ is a minus witness for the continuation.

It remains to check every claim after the overrides.  Let $y$
claim $i$.  Lemma~\ref{lem:label-down} gives $i\notin I_-$.
If $i\in I_+$, it receives the plus override.  If $i\notin I_+$,
Equation~\eqref{eq:label-plus-survives} applies to $y$, the
continuation labels $i$ by $+$, and no override changes it.
Now let $z$ claim $i$.  Lemma~\ref{lem:label-up} gives
$i\notin I_+$.  If $i\in I_-$, it receives the minus override.
Otherwise Equation~\eqref{eq:label-minus-survives} applies to $z$,
and the continuation's minus label is unchanged.  In particular
$I_-\cap I_+$ contains no claim of either sign, so the stated
priority between overrides causes no conflict.  This completes
the induction.
\end{proof}
